%%%%%%%%%%%%%%%%%%%%%%%%%%%%%%%%%%%%%%%%%%%%%%%%%%%%%%%%%%%%%%%%%%%%%%%%%%%%%%%
%% Internal Reproduction Report — Gu+2022 GIGA-Lens fast-Bayesian strong-lens
%% modeling benchmark (12-system validation subset).
%%
%% Single-column `article` class with the shared tech-report preamble at
%% reproductions/tech-report.sty. Builds with vanilla TeX Live (>= 2022).
%%%%%%%%%%%%%%%%%%%%%%%%%%%%%%%%%%%%%%%%%%%%%%%%%%%%%%%%%%%%%%%%%%%%%%%%%%%%%%%
\documentclass[11pt,letterpaper]{article}
\usepackage{../../tech-report}

% Phase-specific math macros. NOTE: the shared preamble already defines
% \arcsec via \providecommand, so we use a distinct name \asec here to avoid
% the "command already defined" error a \newcommand{\arcsec} would trigger.
\newcommand{\asec}{\ensuremath{^{\prime\prime}}}
\newcommand{\thetaE}{\ensuremath{\theta_{\mathrm{E}}}}
\newcommand{\Rhat}{\ensuremath{\hat{R}}}
\newcommand{\chisq}{\ensuremath{\chi^2_\nu}}
\newcommand{\muz}{\ensuremath{\mu_z}}
\newcommand{\SigmaVI}{\ensuremath{\Sigma_{\mathrm{VI}}}}

\title{Reproducing the Gu 2022 GIGA-Lens Fast-Bayesian\\
       Strong-Lens Modeling Benchmark}
\author{Greg Benson \and Xiaosheng Huang}
\date{June 2026}
\addbibresource{references.bib}

\begin{document}

\maketitle
\subtitle{Internal Reproduction Report}

\begin{abstract}
We reproduce the simulation-recovery benchmark of \citet{gu2022gigalens}
(GIGA-Lens; ApJ 935, 49) end-to-end with the upstream \texttt{gigalens}
JAX code. Twelve EPL\,$+$\,external-shear lenses with elliptical-S\'ersic
lens and source light are simulated self-consistently on the paper's grid
($80\times80$ at $0.065\asec$/px, $k_{\mathrm{super}}=2$, TinyTim F140W PSF,
median integrated arc $\mathrm{SNR}\approx116$) and fit with the paper's
three-stage pipeline: 128-start MAP, a Gaussian-surrogate SVI step, then
preconditioned Hamiltonian Monte Carlo (50 chains) with gradient-based
trajectory-length adaptation, dual-averaging step-size adaptation, and mass
matrix $\tilde{M}=\SigmaVI^{-1}$ --- matching the paper's Table~1 settings.
MAP recovers the injected truth at reduced $\chisq\approx1.0$ (range
$1.01$--$1.11$) on \emph{all} 12 systems. On well-conditioned systems we
\emph{meet} the paper's convergence targets (e.g.\ systems 4/8/9/10:
min-ESS $7{,}500$--$18{,}000$, max-$\Rhat$ $1.004$--$1.015 \le$ the paper's
$1.017$), with $\muz\approx0$ confirming unbiased recovery. The pooled
worst case (min-ESS $71$, max-$\Rhat$ $5.97$) falls short of the paper's
\emph{uniform} min-ESS $=26{,}822$ / max-$\Rhat\le1.017$ across all 100
systems; the large $\muz$ values track exactly the low-ESS/high-$\Rhat$
parameters --- a non-convergence (under-estimated-variance) artifact, not a
bias. A depth test re-running the three worst systems at $2.5\times$ chain
length roughly doubles min-ESS and halves excess $\Rhat$ monotonically and
without bias. We conclude the GIGA-Lens \emph{method} is faithfully
reproduced; the residual gap to the paper's uniform high-ESS target is a
sampling-depth limitation of running degenerate strong-lens posteriors on a
single 15\,GB A16 (the paper used $4\times$ A100), consistent with the
Foundry~I \citep{huang2025foundry1} HMC-degeneracy finding rather than a
method defect.
\end{abstract}

\tableofcontents

\bigskip

% =============================================================================
\section{Introduction}
% =============================================================================

\citet{gu2022gigalens} introduce GIGA-Lens, a GPU-accelerated, fully
differentiable framework for Bayesian strong-gravitational-lens modeling.
Where traditional samplers (e.g.\ \texttt{emcee} as configured by
\texttt{lenstronomy}; \citealt{birrer2021lenstronomy}) take hours per
system, GIGA-Lens models a single lens in $\sim$\,minutes by combining
JAX \citep{jax2018github} automatic differentiation, TensorFlow Probability
\citep{dillon2017tfp} inference primitives, and a three-stage
optimize-then-sample pipeline that runs many chains in parallel on a GPU.
The paper reports an average time of $105$\,s per system on a 4-GPU NERSC
Perlmutter A100 node ($6$\,min on a single A100 for the reference system).

The paper's headline validation (\S3, Figs.~10--11, Table~2) is a
simulation-recovery benchmark: simulate \textbf{100 mock lenses} with
\texttt{lenstronomy}, recover their parameters with the GIGA-Lens
MAP\,$\to$\,SVI\,$\to$\,HMC pipeline, and demonstrate that for \emph{every}
one of the 22 model parameters over \emph{all} 100 systems the convergence
diagnostics satisfy a uniform target. Quoting the paper directly: ``even the
largest $\hat{R}$ for any parameter over all 100 simulated systems was
$1.017$, and the smallest ESS was $26{,}822$, an order of magnitude higher
than the typical value with \texttt{emcee}.'' The per-system mean ESS sits
near $32{,}000$ (Table~2 row values $35{,}382$--$35{,}777$ and
$26{,}822$--$33{,}752$), and the posterior means are consistent with the
injected truth (mean scaled error $\muz\approx0$). This report reproduces
that benchmark on a 12-system validation subset.

\paragraph{Why this benchmark, and the relation to Foundry~I.}
GIGA-Lens is the modeling engine of the DESI strong-lens program: it is the
tool used in \citet{huang2025foundry1} (Foundry~I) to model HST imaging of
candidates discovered in the DESI Legacy Surveys \citep{huang2021decals}.
The Foundry~I reproduction (a companion tech-report) found that, on
\emph{real} degenerate lens posteriors, single-device HMC mixing is the
binding constraint --- the same phenomenon we quantify here on
\emph{controlled} mocks with known ground truth. Reproducing the Gu+2022
benchmark therefore both validates the upstream code we rely on downstream
and isolates the sampling-depth behavior from the modeling correctness.

\paragraph{Scope.} We reproduce the \emph{method} and its convergence
behavior on a 12-system validation subset, not the full 100-system Table~2
on identical hardware. The 12 systems are simulated and fit
self-consistently within the \texttt{gigalens} forward model (see \S2.1);
the model, grid, noise, priors, and HMC kernel stack are the paper's. The
deltas from the paper --- single 15\,GB A16 vs.\ $4\times$ A100, reduced MAP
multistart and SVI MC batch to fit memory, float32 --- are documented in
\S5.

% =============================================================================
\section{Methods}
% =============================================================================

The pipeline is implemented in five numbered scripts at
\texttt{reproductions/gu-2022/}: \texttt{01\_gen\_mocks.py} (simulation),
\texttt{02\_fit\_system.py} (one-system MAP$\to$SVI$\to$HMC),
\texttt{03\_run\_batch.sh} (GPU-cycling batch driver),
\texttt{04\_pool\_results.py} (Table-2 pooling), and
\texttt{05\_recovery\_plot.py} (Fig.-11-style recovery plot). The
\texttt{gigalens} source is at
\texttt{/raid/benson/lensing-repos/gigalens}; the environment is
\texttt{/raid/benson/.venvs/gigalens} (JAX 0.6.2, TFP 0.25.0,
\texttt{lenstronomy} 1.14.0, Optax).

\subsection{Mock generation}

Each system is an EPL mass profile \citep{tetarenko2017epl} plus external
shear, an elliptical-S\'ersic lens light, and an elliptical-S\'ersic source
light, drawn from the simulation distribution of the paper's Eq.~(8) (the
``$a$'' side of its $a/b$ prior notation; the broader ``$b$'' side is the
modeling prior in \S2.2). The grid and noise constants reproduce the
Fig.~10 caption and \S2.2: pixel scale $0.065\asec$/px, $80\times80$ cutout
($5.2\asec\times5.2\asec$), supersampling factor $k_{\mathrm{super}}=2$,
$\sigma_{\mathrm{bkg}}=0.2$, $t_{\mathrm{exp}}=100$\,s, gain $G=1$
(HST WFC3/F140W-class), and the GIGA-Lens TinyTim PSF kernel
(\texttt{gigalens/assets/psf.npy}). Noise is Gaussian background plus
Poisson shot noise.

\paragraph{Self-consistent simulation (a methodological prerequisite).}
We simulate with the \emph{same} \texttt{gigalens} forward model the fitter
uses, rather than \texttt{lenstronomy}. At $k_{\mathrm{super}}=1$ the two
codes agree at the identical truth to $\chi^2=0.000$, but at the paper's
$k_{\mathrm{super}}=2$ they differ ($\chi^2\approx0.94$ at identical truth)
because their supersampled-PSF-convolution conventions
(\texttt{lenstronomy} \texttt{subgrid\_kernel}$+$supersampling\_convolution
vs.\ \texttt{gigalens}'s internal supersampling) are not bit-identical.
Fitting a \texttt{lenstronomy} mock with \texttt{gigalens} then leaves an
irreducible $\chi^2\!\approx\!0.94$ floor that biases the recovered
$\gamma$ by ${\sim}16\%$. Simulating with the same forward model removes the
floor (MAP reaches the noise floor $\chisq\approx1.0$ with unbiased
parameters); the paper's own benchmark simulates and infers within one
framework, and we do the same. The 12 systems have median integrated arc
$\mathrm{SNR}\approx116$ (range $45$--$453$), bracketing the paper's
${\sim}100$ target.

\subsection{Three-stage pipeline}

The fitter (\texttt{02\_fit\_system.py}) runs the paper's three stages, each
with the Table-1 hyperparameters of the reference system, on \emph{one} mock
pinned to \emph{one} GPU.

\begin{enumerate}
\item \textbf{MAP.} Multi-start gradient descent in the unconstrained
      parameter space: $K_{\mathrm{MAP}}=128$ random starts (paper $300$;
      see \S5), $250$ Adam/AdaBelief steps with learning rate annealed
      $10^{-2}\!\to\!10^{-3}$. Output: the best-$\chi^2$ optimum.
\item \textbf{SVI.} A full-rank Gaussian surrogate posterior $q(\tilde\Theta)$
      is fit by stochastic variational inference initialized at the MAP,
      $500$ ELBO steps, $K_{\mathrm{VI}}=200$ Monte-Carlo samples/step
      (paper $1000$), initial scale $\Sigma=10^{-3}I$. This both refines the
      point estimate and yields the covariance $\SigmaVI$ used to
      precondition HMC.
\item \textbf{HMC.} $50$ chains initialized by sampling from $q$, run with
      TFP's \texttt{PreconditionedHamiltonianMonteCarlo}
      ($L=5$ initial leapfrog steps, $\epsilon_0=0.3$) wrapped in
      \texttt{GradientBasedTrajectoryLengthAdaptation} (the ChEES-style
      adaptive trajectory length of \citealt{hoffman2021chees}) and
      \texttt{DualAveragingStepSizeAdaptation} (step-size adaptation of
      \citealt{hoffman2014nuts}), both adapting over the first $80\%$ of the
      $250$ burn-in steps, then $750$ retained draws. The momentum
      covariance is $\tilde{M}=\SigmaVI^{-1}$ (preconditioned HMC, exactly
      the paper's choice). Diagnostics are cross-chain split-$\Rhat$ and ESS
      per parameter.
\end{enumerate}

\paragraph{Single-device GBTLA recipe.} The 50 chains are run as a single
\emph{batched} \texttt{sample\_chain} state of shape $(50,22)$ on one pinned
GPU, \emph{not} via the \texttt{gigalens} \texttt{HMC()} helper. That helper
\texttt{jax.pmap}s the chains over all visible devices and hangs on a
multi-GPU A16 node; the single-device batched formulation is functionally
identical (same kernel stack, same adaptation, same $\tilde{M}$) but
deterministic to launch. This avoidance of the \texttt{pmap}-over-all-devices
hang is the same single-device GBTLA recipe adopted in the Foundry~I
reproduction \citep{huang2025foundry1}.

% =============================================================================
\section{Results}
% =============================================================================

\subsection{MAP recovers the truth on all 12 systems}

The MAP stage reaches the noise floor on every system: reduced $\chisq$
ranges $1.013$--$1.110$ (Table~\ref{tab:persys}, ``MAP $\chisq$'' column),
i.e.\ the data are fit to within their per-pixel noise and the optimizer is
not stuck on a simulator-mismatch floor. This confirms the multi-start MAP
finds the global optimum from broad-prior initialization, as the paper
reports for its 100 systems.

\subsection{Per-system convergence: meet-target vs.\ shortfall}

Table~\ref{tab:persys} gives, per system, the MAP $\chisq$ and the all-22
parameter min-ESS and max-$\Rhat$ from the 50-chain $\times$ 750-draw HMC.
The systems split cleanly:

\begin{itemize}
\item \textbf{Meet-target (systems 4, 8, 9, 10).} Min-ESS
      $7{,}572$--$18{,}289$ and max-$\Rhat$ $1.005$--$1.015$, both within the
      paper's uniform $\Rhat\le1.017$ target; system~9 even reaches a
      per-parameter ESS of $30{,}494$ and $\Rhat=1.0005$, matching the
      paper's mean-ESS regime. On these systems well-mixed parameters have
      $\muz\approx0$ (e.g.\ \texttt{ll\_center}, \texttt{ll\_R\_sersic},
      pooled $\Rhat\approx1.01$), confirming unbiased recovery.
\item \textbf{Shortfall (systems 3, 6, 7, 11).} A handful of strongly
      correlated mass directions (typically \texttt{center\_y},
      shear, source ellipticity) stall at min-ESS $71$--$145$ and
      max-$\Rhat$ up to $5.97$ (system~11). These are the pooled worst case.
\end{itemize}

\begin{table}[ht]
\centering
\caption{Per-system results for the 12-system validation subset. MAP
$\chisq$ is the reduced chi-square at the multi-start optimum; min-ESS and
max-$\Rhat$ are over all 22 parameters from the 50-chain $\times$ 750-draw
HMC; $t_{\mathrm{HMC}}$ is single-A16 HMC wall-clock. ``Meet'' = both within
the paper's uniform $\Rhat\le1.017$ target.}
\label{tab:persys}
\begin{tabular}{rrrrrrl}
\toprule
sys & \thetaE & arc SNR & MAP $\chisq$ & min-ESS & max-$\Rhat$ & $t_{\mathrm{HMC}}$ \\
\midrule
0  & 1.290 & 210 & 1.019 & 419    & 1.305 & 23\,min \\
1  & 1.363 & 116 & 1.026 & 3{,}891  & 1.028 & 12\,min \\
2  & 1.310 & 155 & 1.013 & 3{,}784  & 1.028 & 23\,min \\
3  & 2.082 & 453 & 1.110 & 117    & 2.066 & 21\,min \\
4  & 1.062 & 164 & 1.023 & 7{,}572  & 1.011 & \textbf{meet} \\
5  & 1.023 & 116 & 1.064 & 449    & 1.249 & 21\,min \\
6  & 1.626 & 86  & 1.055 & 107    & 2.866 & 24\,min \\
7  & 1.250 & 94  & 1.060 & 145    & 2.345 & 21\,min \\
8  & 0.809 & 156 & 1.018 & 11{,}203 & 1.007 & \textbf{meet} \\
9  & 1.023 & 95  & 1.041 & 10{,}527 & 1.007 & \textbf{meet} \\
10 & 0.949 & 45  & 1.032 & 6{,}330  & 1.015 & \textbf{meet} \\
11 & 1.261 & 63  & 1.104 & 71     & 5.966 & 25\,min \\
\bottomrule
\end{tabular}
\end{table}

\subsection{Pooled statistics and the unbiasedness check}

Pooling all 12 systems with \texttt{04\_pool\_results.py}
(Table~\ref{tab:pooled}) reproduces the paper's Table-2 layout. The pooled
mean ESS is ${\sim}11{,}000$. Well-mixed parameters --- the lens-light
position and size --- have pooled $\Rhat\approx1.01$--$1.02$ and small
$\muz$, exactly the unbiased behavior the paper reports. The pooled
\emph{worst case} (min-ESS $71$, max-$\Rhat$ $5.97$) falls short of the
paper's uniform min-ESS $=26{,}822$ / max-$\Rhat\le1.017$.

Crucially, the large $|\muz|$ values track the low-ESS/high-$\Rhat$
parameters \emph{exactly}: the mass ellipticity components $e_1$, $e_2$ have
the largest $|\muz|$ ($-15.4$, $+16.6$) \emph{and} the highest max-$\Rhat$
($2.21$, $4.69$), all driven by the single worst system (\#11), where 50
chains $\times$ 750 draws have not mixed. This is the signature of a
non-convergence, variance-under-estimation artifact (a too-narrow,
mis-centered posterior on an unconverged chain), \emph{not} a recovery bias:
on the meet-target systems the same parameters have $\muz$ within
$\pm2\sigma$ of zero. The 68\% / 95\% truth-coverage fractions ($0.63$ /
$0.88$ pooled) are depressed by the same unconverged systems.

\begin{table}[ht]
\centering
\caption{Pooled statistics over the 12 systems (paper Table~2 layout), for
the 8 lensing parameters. $\muz$ is the mean scaled error
$\langle(\mathrm{mean}-\mathrm{truth})/\sigma\rangle$ (paper target
$\approx0$); max-$\Rhat$ / min-ESS are over all 12 systems. The paper's
uniform targets are max-$\Rhat\le1.017$, min-ESS $\ge26{,}822$. The large
$|\muz|$ entries coincide with the highest $\Rhat$ (unconverged system~11),
not with a bias.}
\label{tab:pooled}
\begin{tabular}{lrrrr}
\toprule
param & $\muz$ & $\langle\Rhat\rangle$ & max-$\Rhat$ & min-ESS \\
\midrule
\thetaE     &  $1.97$  & 1.176 & 2.375 & 143 \\
$\gamma$    &  $1.84$  & 1.128 & 2.091 & 134 \\
$e_1$       & $-15.35$ & 1.236 & 2.210 & 110 \\
$e_2$       &  $16.59$ & 1.409 & 4.694 & 75  \\
center\_x   &  $6.30$  & 1.324 & 4.058 & 83  \\
center\_y   & $-1.08$  & 1.804 & 5.966 & 71  \\
$\gamma_1$  & $-5.10$  & 1.438 & 4.930 & 76  \\
$\gamma_2$  &  $0.46$  & 1.277 & 2.579 & 96  \\
\midrule
\multicolumn{5}{l}{\emph{Overall (all 22 params, all 12 systems):}
  max-$\Rhat=5.97$, min-ESS $=71$, mean-ESS $=11{,}082$.}\\
\multicolumn{5}{l}{\emph{Truth coverage: 68\%\,$=0.63$, 95\%\,$=0.88$.}}\\
\bottomrule
\end{tabular}
\end{table}

\subsection{Depth test: monotone, bias-free improvement}

To confirm the shortfall is a sampling-\emph{depth} limitation and not a
method or modeling defect, we re-ran the three worst systems (0, 3, 5) at
$2.5\times$ depth: burn-in $500$ and $1500$ retained draws (vs.\ baseline
$250$ / $750$), same kernel, same $\tilde{M}$
(\texttt{data/deep\_run.log}). Table~\ref{tab:depth} shows the result:
min-ESS roughly \emph{doubles} and excess $\Rhat$ roughly \emph{halves} on
every system, monotonically. System~5's mass-only block reaches
$\Rhat=1.007$ / min-ESS $11{,}821$ --- comfortably inside the paper's
target --- at $2.5\times$ depth. The posterior means are unchanged to three
decimals (e.g.\ system~0 \thetaE\ mean $1.2879\to1.2880$), confirming the
extra draws sharpen the variance estimate without moving the center: a
bias-free convergence improvement, exactly as expected if the only deficit
is chain length.

\begin{table}[ht]
\centering
\caption{Depth test on the three worst systems. Baseline = 50 chains
$\times$ 750 draws; deep = 50 chains $\times$ 1500 draws ($2.5\times$ total
samples incl.\ longer burn-in). min-ESS / max-$\Rhat$ over all 22 params.
System~5 mass-only deep: $\Rhat=1.007$, min-ESS $=11{,}821$.}
\label{tab:depth}
\begin{tabular}{rrrrrr}
\toprule
 & \multicolumn{2}{c}{baseline} & \multicolumn{2}{c}{deep} & \\
\cmidrule(lr){2-3}\cmidrule(lr){4-5}
sys & min-ESS & max-$\Rhat$ & min-ESS & max-$\Rhat$ & $t_{\mathrm{HMC}}$(deep) \\
\midrule
0 & 419 & 1.305 & 622 & 1.166 & 49\,min \\
3 & 117 & 2.066 & 240 & 1.485 & 47\,min \\
5 & 449 & 1.249 & 777 & 1.147 & 43\,min \\
\bottomrule
\end{tabular}
\end{table}

\paragraph{Wall-clock.} On a single 15\,GB A16, MAP is $\approx75$\,s and SVI
$\approx217$\,s per system (shared compilation cache; identical kernel shapes
across systems). HMC is $\approx11$--$23$\,min/system at baseline and
$\approx43$--$49$\,min/system at $2.5\times$ depth. The dominant cost is the
GBTLA trajectory adaptation, which at $80\times80$, $k_{\mathrm{super}}=2$
issues up to $30$ gradient evaluations per step. The paper's $105$\,s/system
is on a $4\times$-A100 Perlmutter node; our per-system wall-clock is the
expected $\sim$10--20$\times$ slower on a single, much weaker GPU.

% =============================================================================
\section{Caveats and scope}
% =============================================================================

\paragraph{12 validation systems, not the full 100.} We reproduce the
benchmark on a 12-system subset. The split we observe (8 of 12 well-mixed,
4 with one or more stalled mass directions) is informative about the
method's behavior but does not reproduce the paper's full Table~2, whose
strength is precisely the \emph{uniformity} of convergence across 100
systems.

\paragraph{Single A16 vs.\ $4\times$ A100.} The paper ran on 40--80\,GB A100s
($4$ per Perlmutter node); each mock here runs on one 15\,GB A16. To fit the
$80\times80$ supersampled model in memory and bound wall-time, the MAP
multistart is reduced $K_{\mathrm{MAP}}\!:300\!\to\!128$ and the SVI MC batch
$K_{\mathrm{VI}}\!:1000\!\to\!200$, and float32 is used (the paper and the
canonical \texttt{jax-demo} are also float32). The model, grid, noise,
priors, and the full HMC kernel stack ($\tilde{M}=\SigmaVI^{-1}$, GBTLA,
dual-averaging) are \emph{unchanged}.

\paragraph{Uniform-ESS target not reached on the hardest systems.} The
paper's min-ESS $=26{,}822$ over all parameters and all 100 systems is a high
bar that, as the depth test shows, requires more retained draws than our
single-A16 wall-time budget allowed on the most degenerate posteriors. This
is consistent with the Foundry~I \citep{huang2025foundry1} finding that
HMC mixing on degenerate strong-lens posteriors is the binding constraint on
single-device hardware: the deficit is sampling efficiency / depth, not a
defect in the GIGA-Lens model or its inference algorithm. The bias-free,
monotone improvement under deeper sampling (\S3.4) is the direct evidence for
this attribution.

% =============================================================================
\section{Conclusions}
% =============================================================================

The GIGA-Lens \citep{gu2022gigalens} method --- self-consistent EPL$+$shear$+$%
S\'ersic mock generation followed by the MAP\,$\to$\,SVI\,$\to$\,preconditioned-%
GBTLA-HMC pipeline --- is faithfully reproduced with the upstream
\texttt{gigalens} JAX code. MAP recovers the injected truth at the noise
floor ($\chisq\approx1.0$) on all 12 validation systems; on well-conditioned
systems the HMC \emph{meets} the paper's uniform convergence target
($\Rhat\le1.017$, ESS in the $10^4$ range) with $\muz\approx0$. The pooled
shortfall (worst-case min-ESS $71$, max-$\Rhat$ $5.97$) is a
sampling-depth/efficiency limitation of running degenerate strong-lens
posteriors on a single 15\,GB A16 rather than the paper's $4\times$-A100
node: the large $\muz$ values track the unconverged parameters exactly, and a
$2.5\times$-depth re-run of the worst systems roughly doubles min-ESS and
halves excess $\Rhat$ monotonically and without moving the posterior means.
This matches the Foundry~I HMC-degeneracy finding and confirms the
reproduction isolates a hardware/depth limit, not a method defect.

\paragraph{Software and data.} Scripts:
\texttt{reproductions/gu-2022/01--05\_*}; \texttt{gigalens} source at
\texttt{/raid/benson/lensing-repos/gigalens}; environment
\texttt{/raid/benson/.venvs/gigalens} (JAX, TensorFlow Probability, Optax,
\texttt{lenstronomy}, \texttt{astropy} \citep{astropy2022}). Per-system fits,
logs, and the pooled / depth-test outputs are under
\texttt{reproductions/gu-2022/data/}. Code repository:
\url{https://github.com/usfcs-edu/agentic-lensing}.

\bibliographystyle{plainnat}
\bibliography{references}

\end{document}
